\documentclass[a4paper,11pt,oneside]{article}


%%%% Packages: %%%%%%%%%%%%%%%%%%%%%%%%%%%%%%%%%%

\usepackage[utf8]{inputenc}
\usepackage[T1]{fontenc}

\usepackage[english]{babel} 

\usepackage{cite}

\usepackage{geometry}
\geometry{textwidth=16cm,textheight=26cm}

%%%% Preamble: %%%%%%%%%%%%%%%%%%%%%%%%%%%%%%%%%%

\title{D3. Scope and research question: <Title>
}

\author{
    <First name> <Last name>
    \and <First name> <Last name>
    \and <First name> <Last name>}
    
\date{
    <Group X>, INFOB302 - IDS (20<XX> - 20<XX>)
}
    
%%%% Document: %%%%%%%%%%%%%%%%%%%%%%%%%%%%%%%%%%

\begin{document}

\maketitle

% ---------------------------------------
\section{Introduction}
% ---------------------------------------

Introduce the topic of your research. Explain the motivation for this research and the problem that you want to solve. This should be done in a way that is understandable for a non-expert (meaning a reader without domain knowledge) reader. Do not forget to add adequate citations if you are using information from another source.

In particular, do not forget to explain the following elements:
the scope: Give the scope of your research. State its general objective and the topics addressed and not addressed by this research; The objectives: State the specific objectives of your research. Explain the gain that is expected in terms of empirical and theoretical knowledge.

% ---------------------------------------
\section{Background}
% ---------------------------------------

Explain the different elements required to understand what your research is about. This section may contain several subsections (one per element). Do not forget to add adequate citations.

\subsection{Element 1}


\subsection{Element 2}

% ---------------------------------------
\section{Related work}
% ---------------------------------------

Explain how your research is related to the state of the art. This should be done by comparing your research to the most relevant related work. This comparison should be done in terms of the research question(s) and the research method(s) used. 
In particular, specify if your research is relying on a similar approach for a different problem; addressing the same problem using a different approach; addressing the same problem using the same approach but with a different goal; or addressing the same problem using the same approach with the same goal. 

Do not forget to add adequate citations.

% ---------------------------------------
\section{Research question(s)}
% ---------------------------------------

Provide the main research question(s) that the research will answer.

\begin{itemize}
    \item \textbf{RQ.1} ...
\end{itemize}

Provide an explanation and justification for each research question and subdivide it into sub-research-questions or hypothesis: 

\begin{itemize}
    \item \textbf{RQ.1.1} ...
    \item \textbf{RQ.1.2} ...
\end{itemize}

Provide an explanation and justification for each sub-research question. This should state how the sub-research question contributes to answering the global research-question and indications on how it will be answered. For instance, by providing a general description of the kind of research method and analysis that could be applied. 

Do not forget to add adequate citations if you are using information from another source \cite{...}. For instance, it is required in a research question to compare your own work to the state of the art. This requires adequate citation of the original papers describing related approaches, research methods, etc.

\bibliographystyle{plain}
\bibliography{biblio}

\end{document}
